% Level 2, Week 11 answers -- Project Studio 2
\answersweek{11}{Project Studio 2}{Answers and scoring notes for workbook chapter 11 and homework 11.}

\section{Workbook exercises}
\begin{answer}[Exercise 11.1: Add your data (M3)]
Example for an arcade game: list \texttt{stars} (the falling star turtles); dictionary
\texttt{SETTINGS} (speed, step, edges) or \texttt{state} (score, lives, running); file
\texttt{highscore.txt} (the best score). Example for an Art Studio: list \texttt{shapes}; each shape is
a dictionary \texttt{\{"x": …, "y": …, "size": …, "color": …\}}; file \texttt{design.txt}, one shape
per line. \textbf{M3 data met} when the data comes back after closing and re-running the program, and
a missing file does not crash it (a \texttt{try}/\texttt{except FileNotFoundError}, as in week 7).
\end{answer}

\begin{answer}[Exercise 11.2: Write three tests]
The arcade tests in the lesson plan (week 11, \texttt{test\_mygame.py}) pass against the skeleton
from week 10 plus the two file functions. For Art Studio teams, this module and its tests were run
with pytest (3 passed):
\begin{pycode}[artstudio.py -- save and load a design (tested)]
def shape_to_line(shape):
    """Turn one shape dictionary into one line of text for the file."""
    values = [shape["x"], shape["y"], shape["size"], shape["color"]]
    return ",".join(str(v) for v in values) + "\n"

def line_to_shape(line):
    """Turn one line from the file back into a shape dictionary."""
    x, y, size, color = line.strip().split(",")
    return {"x": int(x), "y": int(y), "size": int(size), "color": color}

def save_design(shapes, filename):
    with open(filename, "w") as f:
        for shape in shapes:
            f.write(shape_to_line(shape))

def load_design(filename):
    with open(filename) as f:
        return [line_to_shape(line) for line in f if line.strip()]
\end{pycode}
\begin{pycode}[test\_artstudio.py (tested)]
from artstudio import line_to_shape, load_design, save_design, shape_to_line

def test_shape_becomes_a_line():
    shape = {"x": 10, "y": -20, "size": 30, "color": "#0072B2"}
    assert shape_to_line(shape) == "10,-20,30,#0072B2\n"

def test_line_becomes_a_shape():
    assert line_to_shape("5,6,7,#E69F00\n")["size"] == 7

def test_save_then_load_gives_the_same_design():
    shapes = [{"x": 1, "y": 2, "size": 3, "color": "#009E73"}]
    save_design(shapes, "test_design.txt")
    assert load_design("test_design.txt") == shapes
\end{pycode}
\textbf{Requirement met} when at least 3 tests call the team's own functions, each has an
\texttt{assert}, they all pass, and at least one would fail if the code were broken (ask the team to
show you). \textbf{Not met}: \texttt{assert True}, tests of drawing functions only, or tests that use
the real high-score file.
\end{answer}

\begin{answer}[Exercise 11.3: Playtest feedback and bug list]
A useful bug line reads like: \enquote{I held the right arrow | the basket went off the screen | it
should stop at the edge | MUST}. That one is a logic bug: a test like
\texttt{assert next\_x(280, 1) == 280} shows it, then clamping fixes it. \enquote{The stars are hard to
see on the blue background | SHOULD} is a polish item. Accept any honest list; check that MUST items
were fixed first.
\end{answer}

\begin{answer}[Exercise 11.4: Accessibility review and polish (M4)]
Good changes: different shapes for good and bad objects (\texttt{"circle"} vs \texttt{"square"}),
a score written with \texttt{font=("Arial", 18, "normal")} or larger, an on-screen line explaining the
keys, palette colors with enough contrast against the background, a slower first level. No
flashing more than 3 times per second. \textbf{M4 met} when at least one change was made and the
checklist is honestly checked.
\end{answer}

\section{Homework}
\begin{answer}[Core: three passing tests and one fix]
The last line should read like \texttt{3 passed in 0.05s} (more tests are fine). If it says
\texttt{collected 0 items}, the file or function names do not start with \texttt{test\_}, or the shell
was in the wrong folder (\texttt{ls} shows the files in the current folder).
\end{answer}

\begin{answer}[Challenge: test first]
Any test that failed before the change and passed after is a full answer. Seeing it fail first proves
the test can catch the bug; a test that has never failed might not be checking anything. Example
(also the lesson's stretch challenge), run with pytest (3 passed):
\begin{pycode}[scores.py and test\_scores.py (tested)]
def add_score(scores, new_score):
    """Return a new list with the best five scores, biggest first."""
    return sorted(scores + [new_score], reverse=True)[:5]

# --- test_scores.py ---
from scores import add_score

def test_first_score_in_an_empty_list():
    assert add_score([], 7) == [7]

def test_small_score_does_not_get_in():
    assert add_score([50, 40, 30, 20, 10], 5) == [50, 40, 30, 20, 10]

def test_new_best_score_goes_to_the_top():
    assert add_score([50, 40, 30, 20, 10], 99) == [99, 50, 40, 30, 20]
\end{pycode}
\end{answer}

\begin{answer}[Creative: how to play]
Look for the goal, the keys and the ending in three short lines, for example \enquote{Catch the
stars, dodge the rocks. ← → to move. Three misses and the game is over.} The real test: a family
member can start playing without asking a question.
\end{answer}
